% ATTEMPT_STATUS: unresolved
% TOP_PROBLEM: {"problemId": "problem.nearby-lagrangian-conjecture", "canonicalTitle": "Nearby Lagrangian Conjecture", "releaseRank": 270, "primaryDomain": "geometry_topology", "primaryDomainLabel": "Geometry & topology", "displayStatus": "open", "releaseStatus": "open"}
% !TeX program = lualatex
\documentclass[11pt]{article}
\usepackage[margin=25mm]{geometry}
\usepackage{fontspec}
\setmainfont{Latin Modern Roman}
\usepackage{amsmath,amssymb,mathtools}
\usepackage{unicode-math}
\setmathfont{Latin Modern Math}
\usepackage{xurl,hyperref}
\hypersetup{colorlinks=true,urlcolor=blue}
\setcounter{secnumdepth}{0}
\setlength{\parindent}{0pt}
\setlength{\parskip}{5pt}
\setlength{\emergencystretch}{3em}
\begin{document}
\section*{\texorpdfstring{Nearby Lagrangian Conjecture}{Nearby Lagrangian Conjecture}}\label{270}
\subsection{Definitions and mathematical statement}
For the smooth base manifold $N$ in the standard closed-base cotangent formulation, $T^*N$ has canonical one-form $\lambda$ and symplectic form ${\,\mathrm{d}}\lambda$. A closed embedded submanifold $L$ is Lagrangian when $\dim L=\dim N$ and $({\,\mathrm{d}}\lambda)|_L=0$, and exact when $\lambda|_L={\,\mathrm{d}} f$ for some function $f:L\to{\mathbb{R}}$. The conjecture asserts that there is a Hamiltonian isotopy $\Phi_t$ of $T^*N$ with
\[
 \Phi_1(L)=\text{the zero section }N.
\]
The isotopy is generated by a time-dependent Hamiltonian in the source's support convention. Homotopy equivalence of the projection $L\to N$, or abstract diffeomorphism of $L$ and $N$, is weaker than this ambient Hamiltonian-isotopy conclusion.
\par\smallskip\textit{Formulation and concept references:} \href{https://www.math.stonybrook.edu/~fzheng/seminar-f24/notes/week-1-2.pdf}{[S1]}.

\subsection{Short English statement}
Is every closed exact Lagrangian in a cotangent bundle obtainable from its zero section by a Hamiltonian isotopy?
\subsection{Research attempt}
\paragraph{Scope and source check.}
GPT-6 Astra Ultra prepared this unresolved attempt on 27 September 2026
by reviewing the formulation source and checking explicit symplectic
calculations. The base is closed; for the proof below take a connected
component of the base at a time. The source distinguishes a homotopy
classification from the requested Hamiltonian isotopy. Its surrounding
lecture notes contain some explicitly tentative arguments, so those
arguments are not adopted as theorem inputs here. The elementary route
chosen below solves the case in which the cotangent projection has no
critical points, and identifies the exact additional geometric task in
the general case. All Hamiltonians constructed below can be compactly
supported, which meets the usual closed-base support convention.

\paragraph{Exact graphs admit an explicit isotopy.}
Write $\pi:T^*N\to N$ and use the convention
$\omega=d\lambda=\sum_i dp_i\wedge dq_i$ and
$\iota_{X_H}\omega=-dH$. If $\alpha$ is a one-form on $N$, its graph
$j(q)=(q,\alpha_q)$ satisfies
\[
 j^*\lambda=\alpha,\qquad j^*\omega=d\alpha.
\]
Thus the graph is Lagrangian precisely when $\alpha$ is closed, and
exact precisely when $\alpha=df$ for a globally defined function.
For $H(q,p)=f(q)$ the Hamilton equations in this convention are
$\dot q=0$, $\dot p=-df(q)$. Its time-one map sends the graph of
$df$ to the zero section.

The Hamiltonian just written is not compactly supported in the fibers.
This is easily repaired without changing the prescribed isotopy. The
swept set
\[
 \{(q,s\,df(q)):q\in N,\ 0\le s\le1\}
\]
is compact. Choose a smooth compactly supported cutoff $\chi$ equal
to one on a neighborhood of it. For $\widehat H=\chi f$, both the
Hamiltonian and its differential agree with those of $H$ on that
neighborhood. Uniqueness for the Hamiltonian ordinary differential
equation shows that the points on the moving graph follow exactly the
same paths. The cutoff flow exists throughout the required interval
because its vector field has compact support. This proves the full
ambient conclusion for exact graphs, including the support condition.

\paragraph{A projection with no critical points must be a graph.}
There is a stronger elementary subcase. Suppose $L$ is nonempty, closed,
embedded and exact, and $\pi|_L$ is a local diffeomorphism. Compactness
makes this map proper. Its image is both open and closed in connected
$N$, hence is all of $N$; it is a finite covering. I claim that its
degree is one. This conclusion uses exactness and embeddedness together,
and does not assume that a covering of the base is automatically trivial.

Choose $f:L\to\mathbb R$ with $df=\lambda|_L$. If the covering had
at least two sheets, the space
\[
 C=\{(x,y)\in L\times L:\pi(x)=\pi(y),\ x\ne y\}
\]
would be a nonempty compact smooth manifold. For a finite covering the
diagonal is both open and closed in the fiber product, so removing it
preserves compactness; there is no sequence of distinct sheets converging
to the diagonal. The function $F(x,y)=f(x)-f(y)$ has a maximum on $C$.
At such a maximum, vary the common base point in an arbitrary tangent
direction $v\in T_qN$, lifting it to both sheets. Write
$x=(q,p_x)$ and $y=(q,p_y)$. The differential of $F$ is
\[
 dF(v)=p_x(v)-p_y(v).
\]
It vanishes for every $v$, so $p_x=p_y$. Consequently $x=y$ as points
of the embedded submanifold of $T^*N$, contradicting their membership
in $C$. The cover therefore has degree one and is a diffeomorphism.
Its inverse identifies $L$ with the graph of $d(f\circ(\pi|_L)^{-1})$.
The preceding compactly supported Hamiltonian sends it to $N$.

This proof also excludes a disconnected exact multisection over a
connected base under the same nonsingularity hypothesis: the fiber
product argument does not require connectedness of $L$. It does not
exclude disconnectedness or prove the conjecture for a general exact
Lagrangian whose projection has critical points. The compact covering
structure, and specifically compactness of $C$, was essential.

\paragraph{Why exactness cannot be removed.}
On $T^*\mathbb T^m$, a constant nonzero covector $a$ defines a closed
Lagrangian graph. Its Liouville periods on the coordinate circles are
the components of $a$, with the corresponding circle-length convention.
If one is nonzero, the graph is not exact and cannot be Hamiltonian
isotopic to the zero section. Here is the invariance calculation.
For a Hamiltonian flow $\Phi_t$,
\[
 \frac{d}{dt}\Phi_t^*\lambda
 =\Phi_t^*\bigl(\iota_{X_{H_t}}d\lambda+
                           d(\lambda(X_{H_t}))\bigr)
 =d\,\Phi_t^*\bigl(\lambda(X_{H_t})-H_t\bigr).
\]
Integrating in time shows that $\Phi_t^*\lambda-\lambda$ is exact.
Its integral around every closed loop is zero. Thus a nonzero period
cannot become a zero-section period under such a flow.

The example is a stress test of the hypotheses, not a counterexample
to the problem. It also distinguishes Hamiltonian translation by $df$
from symplectic translation by an arbitrary closed one-form. The latter
preserves $d\lambda$ but need not have a globally defined Hamiltonian.
Ignoring this distinction would incorrectly prove an assertion without
the necessary exactness assumption.

\paragraph{Exact Lagrangian paths would be enough.}
A useful reduction is that a smooth path of closed embedded exact
Lagrangians can be extended, at the level of their images, to a
Hamiltonian isotopy. To check the statement, let
$j_t:L_0\to T^*N$ be a smooth family of Lagrangian embeddings and
$V_t=\partial_tj_t$ its velocity along the image. Define
\[
 \beta_t=j_t^*(\iota_{V_t}\omega).
\]
Differentiating $j_t^*\omega=0$ gives $d\beta_t=0$. If
$j_t^*\lambda=df_t$ with primitives chosen smoothly in $t$, Cartan's
formula gives
\[
 \beta_t=d\bigl(\partial_tf_t-\lambda(V_t)\circ j_t\bigr).
\]
Write the last primitive as $h_t$. Smooth primitives can be obtained
by fixing base values and integrating the smoothly varying exact forms;
there is no changing period obstruction.

Choose a smooth Hamiltonian near the moving image with
$H_t\circ j_t=-h_t$, and extend it with compact support. Such an
extension exists using a tubular neighborhood of the compact family
of embeddings. For every tangent vector $w$ to $L_t$,
\[
 \omega(X_{H_t}-V_t,w)=0.
\]
The symplectic orthogonal of a Lagrangian tangent space is that same
tangent space. Hence $X_{H_t}-V_t$ is tangential: the Hamiltonian flow
and $j_t$ move the same submanifold, differing only by a time-dependent
reparametrization. This proves the reduction. It does not construct the
required path from an arbitrary $L$ to the zero section.

A merely formal deformation of tangent planes, a homotopy of the
projection, or a path through immersed Lagrangians does not meet this
criterion. The path must consist of actual embedded exact Lagrangians.
At a self-intersection there need not be one compatible value of the
local Hamiltonian prescription on the ambient image. This is one reason
that an immersion-level construction cannot simply be declared an
ambient isotopy.

\paragraph{The fiber-contraction shortcut fails at its endpoint.}
Consider $\delta_s(q,p)=(q,sp)$. For $s>0$ it is a diffeomorphism and
$\delta_s^*\omega=s\omega$, so it takes exact Lagrangians to exact
Lagrangians. But it is generally not symplectic and is not itself a
Hamiltonian map. More decisively, at $s=0$ its restriction to $L$ is
just the projection onto the zero section. Critical points of
$\pi|_L$ can destroy the immersion condition, and several points can
have the same image. Thus the family need not be a smooth path of
embeddings including its endpoint. Applying the preceding extension
lemma on $s\in[\epsilon,1]$ for every positive $\epsilon$ does not
supply a limiting Hamiltonian diffeomorphism at zero.

Even uniform convergence of the images as sets would not fix this:
invertibility and derivative bounds need not survive such a limit.
The graph subcase avoids the problem because its projection is already
a diffeomorphism and the explicit Hamiltonian supplies a nonsingular
path all the way to the endpoint.

\paragraph{Verification, first gap, and outcome.}
The sign of the Hamilton equations was checked against the source's
$d\lambda$ convention, the cutoff preserves the entire swept graph,
and the covering proof explicitly checks compactness after removing the
diagonal. The period obstruction tests the necessity of exactness.
The first missing step is an embedded exact deformation that removes
the critical points of the cotangent projection, or another construction
of the required Hamiltonian isotopy without doing so. The projection
subcase and path-extension reduction are proved; neither supplies that
step for arbitrary $L$. The nearby Lagrangian conjecture remains
\textbf{UNRESOLVED}, and no novelty is claimed for these elementary
restricted deductions.

\subsection{Sources}
\begin{itemize}
\item[S1]Formal statement, Conjecture 1.5.
\url{https://www.math.stonybrook.edu/~fzheng/seminar-f24/notes/week-1-2.pdf}.
\textit{Formulation source.}
\end{itemize}
\end{document}
